% Source: https://artofproblemsolving.com/community/c3h2791549_2011_math_hour_olympiad__university_of_washington__grades_67
% Stable identity: aops-topic-2791549-q01
The ratio must be juice : acid : water $=1:2:3$, the overall ratio. Call the vials $J_1,A_2,W_3$ (initially full) and $E_1,E_2,E_3$ (empty), subscripts giving capacities, and write contents as (juice, acid, water).
\begin{enumerate}
\item $A_2\to E_1$: $E_1=(0,1,0)$, $A_2=(0,1,0)$.
\item $J_1\to A_2$: $A_2=(1,1,0)$, full.
\item $A_2\to J_1$: halves it; $J_1=A_2=(\dfrac12,\dfrac12,0)$.
\item $E_1\to E_2$, then $W_3\to E_2$ until full: $E_2=(0,1,1)$, $W_3$ holds $2$ oz of water.
\item $E_2\to E_1$: halves it; $E_1=E_2=(0,\dfrac12,\dfrac12)$.
\item $J_1\to E_3$ and $E_1\to E_3$: $E_3=(\dfrac12,1,\dfrac12)$.
\item $W_3\to E_1$ (fills it with water, leaving $1$ oz in $W_3$), then $E_1\to E_3$: $E_3=(\dfrac12,1,\dfrac32)$, full.
\item $A_2\to W_3$ and $E_2\to W_3$: $W_3=(\dfrac12,1,\dfrac32)$, full.
\item $W_3\to J_1$ and then $W_3\to A_2$: these fill $J_1$ and $A_2$ from a uniform mixture.
\end{enumerate}
Now $J_1$, $A_2$ and $E_3$ are full, each with ratio $\boxed{1:2:3}$. Every step is a legal pour.
